% 第 11 讲 映射复形与投射消解
% 源：src/lecture11.txt（264 行）
\section*{第 11 讲\quad 映射复形与投射消解}

本讲始终设 $A$ 是结合代数（可以非交换），例如 $A=\Z$ 或 $\Z[G]$，其中 $G$ 是（通常是有限的）群。回忆我们想要的规则是 $X\mapsto X_{hG},\,X^{hG}$：$X$ 是 $G$-模的复形，而 $X_{hG}$、$X^{hG}$ 是阿贝尔群的复形。这一构造的指导性「公理」将是：

\begin{enumerate}[label=(\arabic*),leftmargin=2.4em]
\item 若 $X$ \textbf{无环}，即对一切 $i\in\Z$ 有 $H_i(X)=0$，则希望 $X_{hG}$ 与 $X^{hG}$ 也无环。
\item $X_{hG}$ 与 $X^{hG}$ 都应与锥交换：若 $f\colon X\to Y$ 是 $G$-模复形的映射，则 $\operatorname{hCoker}(f)^{hG}\simeq\operatorname{hCoker}(X^{hG}\to Y^{hG})$。
\end{enumerate}

第二条条件相对容易满足，因为由断言 10.1 的证明，它只要求与有限直和及平移交换。

\begin{enumerate}[label=(\arabic*),leftmargin=2.4em,start=3]
\item 该构造应与不变子模、余不变商有关。即若 $X=(\cdots\to0\to M\to0\to\cdots)$ 只在次数 $0$ 有项，则 $H^0(X^{hG})=M^G$，且
\[
H^0(X_{hG})=M_G=\Z\otimes_{\Z[G]}M=M\Big/\sum_{g\in G}(g-1)M.
\]
\end{enumerate}

一个朴素而不正确的尝试是定义
\[
X^{hG}:=\Bigl(\cdots\to(X^{-1})^G\xrightarrow{\ d\ }(X^0)^G\xrightarrow{\ d\ }(X^1)^G\to\cdots\Bigr),
\]
其中
\[
X:=(\cdots\to X^{-1}\xrightarrow{\ d\ }X^0\xrightarrow{\ d\ }X^1\to\cdots).
\]
它平凡地满足 (2) 与 (3)；并注意它确实良定义，因为微分与群自同构交换。它还满足 (1) 的一个弱化版本：若 $X\simeq0$（即零复形），则 $X^{hG}$ 与 $X_{hG}$ 也同伦等价于 $0$。然而，这个构造不保持无环复形。具体地，若 $G:=\Z/2\Z$ 通过乘以 $1$ 与 $-1$ 作用在 $\Z$ 上，则
\[
\Bigl(\cdots\to0\to\Z\xrightarrow{\times2}\Z\to\Z/2\to0\to\cdots\Bigr)^{hG}
=\Bigl(\cdots\to0\to0\to0\to\Z/2\to0\to\cdots\Bigr),
\]
而后者不是无环的！

一个更省事的论证是：若该构造保持无环复形，则由 (10.1) 的构造，任意两个无环复形之间映射的锥必无环；又由假设它与锥交换，故由推论 10.11 它也会保持拟同构。但由断言 10.12，对 $G$-模的单射 $M\hookrightarrow N$，有 $\operatorname{hCoker}(i)\xrightarrow{\mathrm{qis}}\operatorname{Coker}(i)=N/M$（此后「qis」均指拟同构）。于是，若朴素不变子模保持无环复形，它就会保持余核，而我们知道这是错的。

注意，对结合代数 $A$，若 $A=\Z[G]$，则 $M\mapsto M^G=\operatorname{Hom}_{\Z[G]}(\Z,M)$，其中 $G$ 以平凡作用作用于 $\Z$。于是，对每个结合代数 $A$ 与 $A$-模 $M$ 都有一类一般的问题。

\textbf{定义 11.1.} 设 $X,Y$ 是 $A$-模的复形。\textbf{映射复形} $\operatorname{Hom}_A(X,Y)$ 是由
\[
\operatorname{Hom}_A^i(X,Y):=\prod_{j\in\Z}\operatorname{Hom}_A(X^j,Y^{j+i}),\qquad d^if:=df-(-1)^ifd
\]
定义的阿贝尔群复形（符号交替以保证微分平方为零）。可图示如下：
\[
\operatorname{Hom}^{-1}(X,Y)\xrightarrow{\ h\mapsto dh+hd\ }\operatorname{Hom}^0(X,Y)\xrightarrow{\ f\mapsto df-fd\ }\operatorname{Hom}^1(X,Y).
\]
其中各项按 $\operatorname{Hom}_A^i(X,Y)=\prod_j\operatorname{Hom}_A(X^j,Y^{j+i})$ 展开为分量：$\operatorname{Hom}^{-1}$ 的分量形如 $h^{-1}\colon X^{-1}\to Y^{-2}$、$h^0\colon X^0\to Y^{-1}$、$h^1\colon X^1\to Y^0$；$\operatorname{Hom}^0$ 的分量形如 $f^{-1}\colon X^{-1}\to Y^{-1}$、$f^0\colon X^0\to Y^0$、$f^1\colon X^1\to Y^1$；$\operatorname{Hom}^1$ 的分量形如 $g^{-1}\colon X^{-1}\to Y^0$、$g^0\colon X^0\to Y^1$、$g^1\colon X^1\to Y^2$。原图把这三列分量并排画出，并用 $h\mapsto dh+hd$、$f\mapsto df-fd$ 标出相邻两列之间的微分；其中 $d$ 分别指 $X$ 与 $Y$ 的微分。

\textbf{断言 11.2.} 对任意 $A$-模复形 $X,Y$，$\operatorname{Hom}_A(X,Y)$ 是复形。

注意，复形映射 $f\colon X\to Y$ 等价于元素 $f=(f^i)\in\operatorname{Hom}^0_A(X,Y)$ 且 $df=0$，其中 $d$ 指 $\operatorname{Hom}(X,Y)$ 上的微分。$f$ 的零伦同样等价于元素 $h\in\operatorname{Hom}^{-1}_A(X,Y)$ 且 $dh=f$，$d$ 同上。于是 $H^0\operatorname{Hom}(X,Y)$ 等价于 $X\to Y$ 的映射模同伦的等价类。因此这一构造推广了同调代数中许多重要概念。

\textbf{例 11.3.} 若 $X:=(\cdots\to0\to A\to0\to\cdots)$，$A$ 位于次数 $0$，则 $\operatorname{Hom}_A(X,Y)=Y$。因此这样的 $X$ 就是我们即将称为投射的对象。

\textbf{定义 11.4.} $A$-模的复形 $P$ 称为\textbf{投射的}（或称同伦投射、$K$-投射等；此概念由 Spaltenstein 定义），若对每个无环的 $A$-模复形 $Y$，$\operatorname{Hom}_A(P,Y)$ 也无环。

上面的困难在于：$\Z$ 作为 $\Z[G]$-模的复形并不投射。我们将说明，在某种意义下可以把 $\Z$「唯一地」替换为投射对象。

\textbf{引理 11.5.} 若 $P$ 是 $A$-模复形，$P^i=0$ 对一切 $i\gg0$（即 $P$ 的非零项在指标上有上界），且每个 $P^i$ 作为 $A$-模是投射的，则 $P$（作为复形）是投射的。

回忆下述定义：

\textbf{定义 11.6.} $A$-模 $P^i$ 作为 $A$-模\textbf{投射}，若以下等价条件之一成立：
\begin{enumerate}[label=(\arabic*),leftmargin=2.4em]
\item $\operatorname{Hom}_A(P^i,-)$ 保持满射；
\item $\operatorname{Hom}_A(P^i,-)$ 是正合函子；
\item $P^i$ 是自由模的直和项；
\item 对任意满射 $N\twoheadrightarrow M$，每个 $A$-模映射 $P^i\to M$ 都可提升为映射 $P^i\to N$，使下图交换：
\[
\begin{array}{ccc}
& & N\\
& \nearrow & \big\downarrow\\
P^i & \longrightarrow & M\longrightarrow0.
\end{array}
\]
\end{enumerate}

我们简要说明这些等价性。(1) 与 (4) 显然等价，因为 (4) 说的正是：若 $N\twoheadrightarrow M$，则 $\operatorname{Hom}_A(P^i,N)\twoheadrightarrow\operatorname{Hom}_A(P^i,M)$。条件 (2) 与 (1) 平凡地等价。为证 (3)，取 $M:=P^i$，$N$ 为某个满射到 $P^i$ 的自由模（例如以 $P^i$ 的全部元素为基，把对应元素送到对应元素）。于是 (4) 给出 $N\twoheadrightarrow P^i$ 的一个分裂，把 $P^i$ 实现为 $N$ 的直和项（原文此处印作 $M$ 的直和项，但此处 $M=P^i$，应为 $N$；已径改）。容易看出投射模的直和项仍投射；至于反向，只需注意自由模是投射的。

\textbf{断言 11.7.} $A$-模 $P$ 作为 $A$-模投射，当且仅当它作为次数 $0$ 的复形投射。

\textit{证明.} 若 $P$ 作为次数 $0$ 的复形投射，取满射 $f\colon N\twoheadrightarrow M$，并作无环复形
\[
X:=(\cdots\to0\to\ker(f)\to N\to M\to0\to\cdots).
\]
则 $\operatorname{Hom}_A(P,X)$ 为
\[
\cdots\to0\to\operatorname{Hom}_A(P,\ker(f))\to\operatorname{Hom}_A(P,N)\to\operatorname{Hom}_A(P,M)\to0\to\cdots,
\]
故 $\operatorname{Hom}_A(P,-)$ 保持满射，按定义 $P$ 作为 $A$-模投射。

反之，若 $X:=(\cdots\to X^{-1}\xrightarrow{d^{-1}}X^0\xrightarrow{d^0}X^1\to\cdots)$ 是无环复形且 $P$ 作为 $A$-模投射，则
\[
\operatorname{Hom}_A(P,X)=(\cdots\to\operatorname{Hom}_A(P,X^{-1})\to\operatorname{Hom}_A(P,X^0)\to\operatorname{Hom}_A(P,X^1)\to\cdots),
\]
它无环：因为若 $X^{i-1}\twoheadrightarrow\ker(X^i\to X^{i+1})$，则
\[
\operatorname{Hom}(P,X^{i-1})\twoheadrightarrow\operatorname{Hom}(P,\ker(X^i\to X^{i+1}))=\ker\bigl(\operatorname{Hom}(P,X^i)\to\operatorname{Hom}(P,X^{i+1})\bigr),
\]
这里用了假设 $\operatorname{Hom}(P,-)$ 正合、从而保持核。故 $\operatorname{Hom}_A(P,X)$ 无环，$P$ 作为次数 $0$ 的复形投射，如所欲证。\hfill$\square$

\textit{证明（引理 11.5）.} 设 $Y$ 是无环的 $A$-模复形。我们需要下述断言：

\textbf{断言 11.8.} 每个映射 $P\to Y$ 都零伦。

\textit{证明.} 设 $f\colon P\to Y$。我们用降序归纳构造 $f$ 的零伦 $h$。基例：对一切 $i\gg0$（其含义见引理陈述），有 $P^i=0$，故 $f^i=0$，于是可取 $h^i:=0$。现设已对一切 $j>i$ 有映射 $h\colon P^j\to Y^{j-1}$，使下图交换：
\[
\begin{array}{ccccccccc}
\cdots & \longrightarrow & P^{i-1} & \longrightarrow & P^i & \longrightarrow & P^{i+1} & \longrightarrow & \cdots\\
& & {\scriptstyle f}\big\downarrow & & {\scriptstyle f}\big\downarrow & & {\scriptstyle f}\big\downarrow & & \\
\cdots & \longrightarrow & Y^{i-1} & \longrightarrow & Y^i & \longrightarrow & Y^{i+1} & \longrightarrow & \cdots
\end{array}
\]
（图中标 $?$ 的正是待构造的 $h\colon P^i\to Y^{i-1}$）。我们希望构造 $h\colon P^i\to Y^{i-1}$ 使 $dh+hd=f$。注意，对一切 $x\in P^i$，
\[
(d(f-hd))(x)=(df-dhd)(x)=\bigl(df-(f-hd)d\bigr)(x)=(df-fd)(x)=0,
\]
这里中间一步用了归纳假设 $dh=f-hd$（作用在 $P^{i+1}$ 上）。由于 $(f-hd)\colon P^i\to\ker(d\colon Y^i\to Y^{i+1})$，且 $d\colon Y^{i-1}\to\ker(d\colon Y^i\to Y^{i+1})$ 是满射，而 $P^i$ 按假设投射，映射 $f-hd$ 可提升为映射 $h\colon P^i\to Y^{i-1}$ 使 $dh=f-hd$。于是 $dh+hd=f$，$h$ 给出 $f$ 的零伦，如所欲证。\hfill$\square$

由该断言，$H^0(\operatorname{Hom}_A(P,Y))$ 中只有零元（其元素即 $P\to Y$ 的同伦类）；且对每个 $i$，有
$H^i(\operatorname{Hom}_A(P,Y))=H^0(\operatorname{Hom}_A(P,Y[i]))=0$，因为 $Y[i]$ 也无环。故各阶上同调消没，$\operatorname{Hom}_A(P,Y)$ 无环，从而 $P$ 是投射的，如所欲证。\hfill$\square$

我们的计划大致是：证明每个 $X$ 都拟同构于某个投射复形 $P$，即 $P\xrightarrow{\mathrm{qis}}X$，称为 $X$ 的\textbf{投射消解}。此外，$P$ 在某种稍后定义的意义下是「唯一」或「导出」的。于是得到「修正」版本 $\operatorname{Hom}^{\mathrm{der}}_A(X,Y):=\operatorname{Hom}_A(P,Y)$。取 $A:=\Z[G]$ 并选取某个投射消解 $P\xrightarrow{\mathrm{qis}}\Z$（它相当典范，对有限群甚至更简单，虽非严格唯一，但对上同调没有影响），便可定义 $X^{hG}:=\operatorname{Hom}_G(P,X)$。这将满足我们的全部公理：它与不变子模有关（因为 $P$ 类似于 $\Z$），又保持无环复形（因为 $P$ 投射）！

下述命题足以证明第一点，因为我们关心的复形是次数 $0$ 处的 $\Z$，它显然有上界。

\textbf{命题 11.9.} 设 $X$ 是 $A$-模复形，且 $X$ 有上界，即如前 $X^i=0$ 对一切 $i\gg0$。则存在投射消解 $P\xrightarrow{\mathrm{qis}}X$。

\textit{证明.} 不妨设 $X$ 在指标 $0$ 处有上界。取自由模 $P^0$ 经映射 $\alpha^0$ 满射到 $X^0$（如前，取生成元即可，使核恰为生成元之间的关系）。再取 $P^{-1}$ 为自由模，满射到纤维积 $P^0\times_{X^0}X^{-1}$：
\[
\begin{array}{ccccccc}
P^{-1} & \longrightarrow & P^0 & \longrightarrow & 0 & \longrightarrow & \cdots\\
\big\downarrow & & \big\downarrow{\scriptstyle\alpha^0} & & & & \\
X^{-1} & \longrightarrow & X^0 & \longrightarrow & 0 & \longrightarrow & \cdots
\end{array}
\]
该构造保持上同调，因为
\[
H^0X=X^0/\operatorname{Im}(X^{-1})=P^0/P^{-1}=H^0P,
\]
这是由于 $P^{-1}$ 满射到 $\ker(\alpha^0)$，且它在 $P^0/\ker(\alpha^0)\cong X^0$ 中的像等于 $X^{-1}$（因为 $P^0\twoheadrightarrow X^0$）。又因 $P^{-1}\twoheadrightarrow X^{-1}$，可以迭代此过程，构造出 $X$ 的投射消解 $P$，如所欲证。\hfill$\square$

第二个论断是关于投射消解的唯一性，由下述引理给出：

\textbf{引理 11.10.} 设 $P_1$ 与 $P_2$ 都是 $A$-模复形 $X$ 的投射消解。则存在同伦等价 $\gamma$，使下图在同伦意义下交换，即 $\beta\gamma\simeq\alpha$：
\[
\begin{array}{ccccc}
& P_1 & \xrightarrow{\ \gamma\ } & P_2 & \\
& {\scriptstyle\alpha}\big\downarrow & & \big\downarrow{\scriptstyle\beta} & \\
& & X & &
\end{array}
\]

\textit{证明.} 考虑下图：
\[
\begin{array}{ccccc}
& P_1 & \xrightarrow{\ \gamma\ } & P_2 & \\
& {\scriptstyle\alpha}\big\downarrow & & \big\downarrow{\scriptstyle\beta} & \\
& X & \xrightarrow{\ \delta\ } & \operatorname{hCoker}(\beta) &
\end{array}
\]
由于 $\beta$ 按假设是拟同构，$\operatorname{hCoker}(\beta)$ 由推论 10.11 无环。由断言 11.8，复合 $\delta$ 零伦；故由断言 10.8，存在典范映射
\[
P_1\to\operatorname{hKer}\bigl(X\to\operatorname{hCoker}(\beta)\bigr)\simeq P_2,
\]
且为同伦等价，如所欲证。由对称性，反方向也存在这样的映射，故 $\gamma$ 是同伦等价，而该图在同伦意义下平凡地交换。\hfill$\square$

现在可以问这里的 $\gamma$ 有多唯一。答案如下：

\textbf{断言 11.11.} 所有这样的 $\gamma$ 都同伦。

\textit{证明.} 我们模仿引理 11.10 的证明，把单个态射换成 $\operatorname{Hom}$-复形。我们有映射
\[
\operatorname{Hom}(P_1,P_2)\xrightarrow{\ \beta_*\ }\operatorname{Hom}(P_1,X)\to\operatorname{Hom}(P_1,\operatorname{hCoker}(\beta))=\operatorname{hCoker}(\beta_*),
\]
其中 $\beta_*$ 由与 $\beta$ 复合给出，最后的等同出于形式理由。由于 $P_2$ 投射，最后的复形无环（按定义），故由推论 10.11，映射 $\psi\colon\operatorname{Hom}(P_1,P_2)\to\operatorname{Hom}(P_1,X)$ 是拟同构，从而是映射同伦类上的同构。特别地，
% [原文疑误：原印 $H^0\operatorname{Hom}(P_1,P_2)=H^0\operatorname{Hom}(P_2,X)$；但 $\psi:\operatorname{Hom}(P_1,P_2)\to\operatorname{Hom}(P_1,X)$ 是拟同构，且给定的映射是 $\alpha\in H^0\operatorname{Hom}(P_1,X)$，故右端应为 $P_1$。此处径改。]
\[
H^0\operatorname{Hom}(P_1,P_2)=H^0\operatorname{Hom}(P_1,X)
\]
于是既然 $H^0\operatorname{Hom}(P_1,X)$ 中给定了映射，$H^0\operatorname{Hom}(P_1,P_2)$ 中的诱导映射便在同伦意义下良定义（如定义 11.1 之后的讨论所述）。\hfill$\square$

事实上可以证明，同伦之间的所有同伦也都同伦，依此类推；故在建立唯一性方面，这已是我们能期望的最佳结果。

\textbf{定义 11.12.} 两个 $A$-模链复形 $M$ 与 $N$ 的第 $i$ 个 \textbf{Ext}-群定义为 $\operatorname{Ext}^i_A(M,N):=H^i\operatorname{Hom}(P,N)$，其中 $P$ 是 $M$ 的某个投射消解。正如我们刚才所示，此定义与 $P$ 的选取无关。
